\section{Statistics for measurement}
\label{ch:stats}

LFS is a statistic, its validation is statistics, and most objections you will
face are statistical. This chapter covers exactly what the paper uses.

\subsection{Variance, sums of squares, and decomposition}

The variance of numbers $x_1,\dots,x_n$ is the mean squared distance from
their mean; the \emph{sum of squares} is the same without dividing. The key
fact of analysis of variance (ANOVA) is that on a balanced design the total
sum of squares splits exactly into parts attributable to each factor plus a
residual, because the parts are orthogonal.

For two crossed factors (language $l$, sentence $c$) with one observation per
cell and one coordinate, write the cell value as
\[
x_{lc} = \mu + a_l + b_c + r_{lc},\qquad
a_l = \bar x_{l\cdot} - \bar x,\quad b_c = \bar x_{\cdot c} - \bar x,\quad
r_{lc} = x_{lc} - \bar x_{l\cdot} - \bar x_{\cdot c} + \bar x .
\]
Then $\sum_{lc}(x_{lc}-\bar x)^2 = N\sum_l a_l^2 + L\sum_c b_c^2 +
\sum_{lc} r_{lc}^2$, exactly, because $\sum_l a_l = 0$, $\sum_c b_c = 0$, and
the residuals sum to zero along every row and column, so every cross term
vanishes. Figure~\ref{fig:anova} shows the picture on one coordinate. LFS is
this decomposition with vectors instead of numbers (sum over coordinates)
and with the residual left out of the ratio.

\begin{figure}[h]\centering\includegraphics[width=0.8\linewidth]{fig_stats_anova.png}
\caption{One hidden coordinate on a small grid: language offsets, concept slope, residual scatter.}\label{fig:anova}\end{figure}

\textbf{Degrees of freedom} count independent pieces of information:
$L-1$ for languages, $N-1$ for sentences, $(L-1)(N-1)$ for the residual. A
\emph{mean square} is a sum of squares divided by its degrees of freedom.

\textbf{Expected mean squares and variance components.} If the effects are
random with variances $\sigma^2_L$, $\sigma^2_C$ and noise $\sigma^2_E$, then
$\mathbb E[\mathrm{MS}_L] = \sigma^2_E + N\sigma^2_L$,
$\mathbb E[\mathrm{MS}_C] = \sigma^2_E + L\sigma^2_C$,
$\mathbb E[\mathrm{MS}_E] = \sigma^2_E$. Solving gives noise-corrected
estimates of each component (the method of moments); restricted maximum
likelihood (REML) does the same with a non-negativity constraint. This is
the LFS-VC variant in the paper. The pure-noise value of the direct LFS
estimator follows from the same expectations with $\sigma^2_L =
\sigma^2_C = 0$: $(L-1)/(L+N-2)$.

\begin{exercise}{verify the decomposition and the null by simulation}
\begin{lstlisting}
def two_way(X):                       # X: (L, N, D)
    L, N, D = X.shape; mu = X.mean((0, 1)); ml = X.mean(1); mc = X.mean(0)
    ss_l = N * ((ml - mu)**2).sum(); ss_c = L * ((mc - mu)**2).sum()
    ss_r = ((X - ml[:, None] - mc[None] + mu)**2).sum(); tot = ((X - mu)**2).sum()
    return ss_l, ss_c, ss_r, tot
X = rng.normal(size=(12, 120, 64))
ss_l, ss_c, ss_r, tot = two_way(X)
print(abs(ss_l + ss_c + ss_r - tot))                    # ~1e-9: exact identity
lfs = [two_way(rng.normal(size=(12, 120, 64)))[0] for _ in range(50)]
sim = np.mean([a / (a + two_way(rng.normal(size=(12,120,64)))[1]) for a in lfs])
print(sim, 11 / (11 + 119))                              # both ~0.0846
\end{lstlisting}
Then add a language effect (\texttt{X += rng.normal(0, 0.5, (12,1,64))}) and
watch the share rise; add a concept effect and watch it fall.
\end{exercise}

\subsection{Standardization and invariances}

Z-scoring a coordinate subtracts its mean and divides by its standard
deviation, an affine map. Done jointly over all cells, it equalizes
coordinate variances without touching the language or sentence structure
within a coordinate. LFS is then invariant to uniform rescaling of all
vectors, to any positive rescaling of individual coordinates, and to
rotations of the coordinate basis; it is \emph{not} invariant to the choice
of basis before standardization (whitening changes it), to pooling, or to
the grid size. Every one of those was measured, and each dictates a
reporting rule.

\subsection{Correlation}

Pearson's $r$ measures linear association; Spearman's $\rho$ is Pearson on
ranks, so it ignores monotone transformations and is robust to outliers. The
paper reports Spearman. A correlation of 0.5 explains a quarter of the
variance; 0.9 explains 81 percent. Correlation is symmetric and says nothing
about cause.

\subsection{Confounding, residualization, and the confounder adjustment regression}

If languages with more training data have both higher alignment and higher
accuracy, then the raw correlation between alignment and accuracy is partly
``data size predicting data size''. The remedy is to remove the observables
from both variables and correlate what is left (Figure~\ref{fig:deflation}).

\begin{figure}[h]\centering\includegraphics[width=0.9\linewidth]{fig_stats_deflation.png}
\caption{Synthetic illustration: a shared driver inflates the raw correlation; residualizing on it leaves the genuine incremental relationship.}\label{fig:deflation}\end{figure}

The paper's regression, for each (model, language) cell:
\[
\mathrm{logit}\!\Big(\tfrac{\mathrm{acc}-0.25}{0.75}\Big) = \text{model FE} + \beta_1 \log_{10}(\text{tokens}) + \beta_2\,\text{macro-family} + \beta_3\,\text{script} + \beta_4\,\text{fertility} + \alpha .
\]
Three details to be able to explain. The \emph{logit transform above the
chance floor} maps accuracies in $(0.25, 1)$ onto the real line so that
differences near the floor and near the ceiling are comparable.
\emph{Model fixed effects} (one dummy per model) absorb each model's overall
level, so the comparison is across languages within a model. \emph{Cluster-robust
standard errors by language} account for the fact that the same language
appears once per model. The residual $\alpha$ is the ``transfer residual'';
metric validity is the Spearman correlation between the residualized metric
and $\alpha$. A validity check precedes everything: the token coefficient must be
positive and significant, or the proxy is broken.

\begin{exercise}{residualize and correlate}
Generate \texttt{tokens}, a \texttt{metric} that depends on tokens plus noise,
and an \texttt{acc} that depends on tokens and a little on the metric.
Compute Spearman raw; then regress both on tokens, take residuals, and
compute Spearman again. Vary the ``little'' and watch when the confounder-adjusted
correlation vanishes. This is Figure~\ref{fig:deflation}.
\end{exercise}

\subsection{Uncertainty: bootstrap, clusters, and the unit of generalization}

A confidence interval says how much the estimate would move if you redrew
the sample. The bootstrap estimates that by resampling your own data with
replacement many times and recomputing the statistic. The key choice is
\emph{what to resample}. If you want to generalize to new languages,
resample languages; to new models, resample models; to both, resample both
(crossed). Rows within a model are correlated, so resampling rows overstates
certainty about models (Figure~\ref{fig:bootstrap}). Four models cannot
support a model-level claim because a four-model bootstrap has only a handful
of distinct resamples; nineteen can. This single point is why the R1
expansion was run and why every interval in the paper is reported by unit.

\begin{figure}[h]\centering\includegraphics[width=0.9\linewidth]{fig_stats_bootstrap.png}
\caption{Same synthetic data: resampling rows gives a narrow interval; resampling the four models gives a wide one. The second is the right one for a claim about models.}\label{fig:bootstrap}\end{figure}

\emph{Effective sample size.} When 40 languages are highly correlated within
a model, they carry the information of far fewer independent languages
(about 1.2 to 2 in the project's panels). Nominal $n$ overstates evidence;
cluster bootstraps do not.

\emph{Leave-one-out.} Recompute the statistic dropping one family (or model)
at a time. If every estimate stays positive, no single family drives the
result. The content-transfer result stays between 0.44 and 0.60 under
leave-one-family-out.

\begin{exercise}{cluster bootstrap for Spearman}
\begin{lstlisting}
from scipy.stats import spearmanr
def boot_ci(x, y, groups, B=2000):
    ids = np.unique(groups); out = []
    for _ in range(B):
        pick = rng.choice(ids, len(ids), replace=True)
        idx = np.concatenate([np.where(groups == g)[0] for g in pick])
        out.append(spearmanr(x[idx], y[idx]).statistic)
    return np.nanpercentile(out, [2.5, 97.5])
\end{lstlisting}
Build 4 models $\times$ 34 languages with a model-level effect; compare
\texttt{groups = model} against \texttt{groups = language}.
\end{exercise}

\subsection{Multiple comparisons}

Test twenty hypotheses at $p < 0.05$ and one will ``pass'' by chance.
Benjamini--Hochberg adjusts a family of $p$-values to control the false
discovery rate; the paper applies it over its preregistered comparison set.
The deeper protection is preregistration: fewer, named tests.

\subsection{Measurement theory: validity, reliability, Goodhart}

\begin{description}[leftmargin=3.4cm, style=nextline, itemsep=1pt]
\item[construct validity] does the number measure what it claims? Evidence: the exact decomposition, the null, the invariances, the probe contrast, the fault test suite.
\item[predictive validity] does it relate to behavior it should relate to? Evidence: content transfer yes; pooled benchmarks no, with a measured reason.
\item[reliability and ceilings] two independent measurements of the same thing agree at 0.315 on their residuals, so no predictor of one can exceed about $\sqrt{0.315} = 0.56$. Spearman--Brown: averaging parallel forms raises reliability.
\item[Goodhart's law] a measure that becomes a target stops measuring. The Aim 2 reversal is a textbook case; BLEU as a training objective is the MT analogue.
\item[preregistration] write the prediction, criterion, and analysis in a dated file before seeing the outcome, and report every failure. It is why the collapse blind spot can be called ``predicted before observed''.
\item[minimum detectable effect] the smallest effect a design can find; a criterion without one returns verdicts it has no power to support (discrepancy ledger D7).
\end{description}

\subsection{Linear algebra tools, in one paragraph each}

\textbf{PCA and SVD.} Directions of largest variance; the spectrum tells how
many directions matter. \textbf{Effective rank and participation ratio.}
Scalar summaries of that spectrum; they fell by 80 to 90 percent when concept
rank collapsed while LFS moved little. \textbf{Whitening (ZCA).} Rotate
and rescale so the covariance is the identity; it equalizes the spectrum and,
in the paper, collapsed the measured language share for 13 of 14 models,
which is how we know language identity lives in high-variance directions.
\textbf{Procrustes.} The best rotation aligning one point set to another;
the map M3 of the sequence. \textbf{Ridge regression.} Least squares with a
penalty; the hub test and map M4. \textbf{Gini coefficient.} A
concentration measure from 0 (spread evenly) to 1 (all in one place); the
coordinate-energy Gini fell from 0.855 to 0.542 under random rotations,
which is why no ``language neuron'' claim is made.
